% =====================================================================
%  CHAPTER 3 — MISSING & EXPANDED SECTIONS
%  "Hybrid CNN–Transformer Networks with Clinical Metadata for
%   Skin Lesion Classification"  (Thesis Group T2510619)
%
%  HOW TO USE THIS FILE
%  --------------------
%  These sections are written to drop straight into your existing
%  Chapter 3.  To make Chapter 3 map cleanly onto every rubric criterion
%  (CO2, CO3, CO4, CO10, CO11, CO12), re-organise the chapter as follows.
%  Sections you ALREADY HAVE are marked [KEEP]; everything in this file
%  is marked [NEW] or [EXPAND].
%
%    3.1  Specifications and Requirements
%         3.1.1  Dataset and Data Requirements              [KEEP - your current 3.1]
%         3.1.2  Functional Requirements                    [NEW]
%         3.1.3  Non-Functional Requirements                [NEW]
%         3.1.4  System Requirements (Hardware & Software)  [KEEP - your current 3.2]
%         3.1.5  Applicable Standards and Codes of Practice [NEW]
%    3.2  Societal, Health, Safety, Legal and Cultural Impact   [EXPAND - replaces thin "Impacts"]   -> CO3
%    3.3  Environmental Impact and Sustainability               [NEW]                                -> CO4
%    3.4  Ethical Issues                                        [EXPAND - promotes your Ethics Statement] -> CO12
%    3.5  Design Constraints                                    [KEEP - your current 3.3, renumbered]
%    3.6  Project Management Plan                               [NEW]                                -> CO10
%    3.7  Risk Management                                       [NEW]                                -> CO10
%    3.8  Economic Analysis                                     [NEW]                                -> CO11
%
%  Requires: \usepackage{booktabs}  (you already use booktabs-style rules).
%  Anything in [SQUARE BRACKETS + bold "REPLACE"] must be replaced with a
%  real value from your own logs before submission.
% =====================================================================


% ---------------------------------------------------------------------
% 3.1.2  FUNCTIONAL REQUIREMENTS                                  [NEW]
% ---------------------------------------------------------------------
\subsection{Functional Requirements}
\label{sec:functional-requirements}

The functional requirements specify what the system must do, independent of
how well it does it. They follow directly from the goal of a clinically
usable, metadata-aware classifier and are summarised in
Table~\ref{tab:functional-requirements}.

The system must accept a single dermoscopic image and return a probability
distribution over the eight ISIC-2019 diagnostic categories (MEL, NV, BCC,
AK, BKL, DF, VASC, SCC), rather than a single hard label, so that the output
is a ranked set of possibilities suitable for decision support. It must
optionally accept the patient metadata that accompanies each image --- age,
sex, and the general anatomical site of the lesion --- and fold that context
into the prediction. Crucially, it must continue to function when any of
those metadata fields is absent, representing ``missing'' as an explicit,
learnable signal rather than imputing a fabricated value, because a large
fraction of ISIC-2019 records are incomplete. The system must apply a fixed
preprocessing pipeline (colour constancy and resizing to $224\times224$) to
raw input before inference, and must report performance per class and as
balanced multi-class accuracy (BMA) so that the behaviour on rare, often
malignant, classes is always visible. Finally, the whole training and
evaluation procedure must be reproducible from configuration files.

\begin{table}[h]
\centering
\caption{Functional requirements of the proposed system.}
\label{tab:functional-requirements}
\begin{tabular}{p{0.07\textwidth}p{0.85\textwidth}}
\toprule
ID & Requirement \\
\midrule
FR1 & Accept an RGB dermoscopic image and output a probability distribution over the eight ISIC-2019 classes. \\
FR2 & Optionally accept patient metadata (age, sex, anatomical site) and use it as additional diagnostic context. \\
FR3 & Handle any missing metadata field through a learned ``missing'' representation, without imputation. \\
FR4 & Return a ranked list of candidate diagnoses with associated scores (decision support, not a final diagnosis). \\
FR5 & Apply the full preprocessing pipeline (Shades-of-Gray colour constancy, resize/centre-crop to $224\times224$) to raw input. \\
FR6 & Train, validate, and evaluate from declarative configuration files with a fixed random seed. \\
FR7 & Report per-class recall, BMA, macro-F1, macro-AUC, overall accuracy, and a normalised confusion matrix. \\
\bottomrule
\end{tabular}
\end{table}


% ---------------------------------------------------------------------
% 3.1.3  NON-FUNCTIONAL REQUIREMENTS                             [NEW]
% ---------------------------------------------------------------------
\subsection{Non-Functional Requirements}
\label{sec:nonfunctional-requirements}

The non-functional requirements specify the quality attributes the system
must satisfy --- the constraints on \emph{how} it performs its function.
They are the requirements against which the design in Chapter~4 is judged,
and they are listed in Table~\ref{tab:nonfunctional-requirements}.

The central non-functional requirement is efficiency: the lightweight model
must stay under 6M parameters and 1 GMAC so that it can be deployed on the
modest hardware found in a typical clinic or on a portable device. The
delivered lightweight hybrid meets this comfortably at 3.98M parameters and
0.631 GMACs. Closely related is deployability --- the system must train and
run within a 16\,GB memory ceiling --- and portability, since the
implementation must run on a standard PyTorch stack, including under Windows,
where the data loader is restricted to a single worker process. On the
quality-of-output side, balanced multi-class accuracy is the primary
performance requirement, and the model must remain robust to the dataset's
extreme imbalance without the rare-class collapse that a naively trained
classifier would exhibit. Reproducibility (fixed seed, deterministic
pipeline, regenerable figures) and maintainability (a modular codebase split
into preprocessing, datasets, losses, models, training, and analysis) are
required so that every reported number can be traced to its source. Two
further requirements are ethical in nature: the reliance on BMA is itself a
fairness requirement, forcing the rare malignant classes to count equally,
and the prohibition on imputing missing metadata is an honesty requirement,
ensuring the model is never given a fabricated clinical value.

\begin{table}[h]
\centering
\caption{Non-functional requirements and how the design satisfies them.}
\label{tab:nonfunctional-requirements}
\begin{tabular}{p{0.08\textwidth}p{0.40\textwidth}p{0.42\textwidth}}
\toprule
ID & Requirement & How it is met \\
\midrule
NFR1 & Efficiency / footprint: lightweight model $\le$ 6M parameters and $\le$ 1 GMAC. & Lightweight hybrid: 3.98M params, 0.631 GMACs. \\
NFR2 & Deployability: train and run within a 16\,GB memory budget. & AMP, channels-last layout, gradient accumulation, memory-mapped dataset. \\
NFR3 & Accuracy: balanced multi-class accuracy as the primary metric; beat same-budget baselines. & Hybrid 0.6081 BMA (best in sub-6M budget); DEKAN 0.6438. \\
NFR4 & Robustness to class imbalance: no rare-class collapse. & Class-Balanced Focal Loss, weighted sampler; stable training curves (Ch.~5). \\
NFR5 & Reproducibility: deterministic, seed-fixed, configuration-driven. & Single seed (42); figures regenerated from saved outputs (Appendix~A). \\
NFR6 & Maintainability: modular, importable codebase. & Separate preprocessing / datasets / losses / models / training / analysis modules. \\
NFR7 & Portability: standard framework; runs under Windows with zero-worker loading. & PyTorch + \texttt{timm}; memmap dataset compensates for single-worker constraint. \\
NFR8 & Fairness and honesty: rare malignant classes weighted equally; no fabricated metadata. & BMA metric; explicit learned ``missing'' embeddings. \\
NFR9 & Latency: single-image inference fast enough for interactive use. & Low GMAC count (0.631) keeps per-image inference inexpensive. \\
\bottomrule
\end{tabular}
\end{table}


% ---------------------------------------------------------------------
% 3.1.5  APPLICABLE STANDARDS AND CODES OF PRACTICE              [NEW]
% ---------------------------------------------------------------------
\subsection{Applicable Standards and Codes of Practice}
\label{sec:standards}

This thesis is a research study built on a public, de-identified dataset, and
is not a regulated medical device. It is nonetheless useful to state the
professional standards and codes of practice that would govern the system
were it taken toward clinical deployment, both to scope the work honestly and
to make clear which good practices we already follow.

Several of these standards bear directly on the development process. For the
software itself, \emph{IEC 62304} defines the life-cycle processes for
medical-device software, and \emph{ISO 13485} the quality-management system a
manufacturer of such software would operate under. \emph{ISO 14971} governs
the application of risk management to medical devices, and informs the risk
register in Section~\ref{sec:risk-management}. Because the system processes
patient metadata, \emph{ISO/IEC 27001} (information-security management) and
data-protection principles of the kind codified in the GDPR are relevant to
any deployment that handles identifiable records; in this work the data is
de-identified at source and the metadata is used solely as model input.

For machine-learning practice specifically, the system is intended to align
with the \emph{Good Machine Learning Practice for Medical Device Development}
guiding principles published jointly by the U.S.\ FDA, Health Canada, and the
UK MHRA, and with the regulatory framing of AI as a Software-as-a-Medical-Device
(SaMD). In the European context, the EU AI Act classifies AI used as a safety
component of a medical device as high-risk, which would impose transparency,
data-governance, and human-oversight obligations on a deployed version of this
system --- obligations consistent with the assistive, human-in-the-loop role
we adopt throughout. Where clinical imaging systems are involved, the
\emph{DICOM} standard governs image interchange.

Finally, for reporting and reproducibility we follow the spirit of the
established checklists for medical-imaging AI --- in particular CLAIM (the
Checklist for Artificial Intelligence in Medical Imaging) and the
TRIPOD\,+\,AI reporting guideline for prediction models --- by reporting the
full data pipeline, the train/validation/test split strategy, the metric
definitions, per-class results, and the limitations of single-seed,
single-dataset evaluation. The data is used in accordance with the ISIC
data-use terms.


% ---------------------------------------------------------------------
% 3.2  SOCIETAL, HEALTH, SAFETY, LEGAL AND CULTURAL IMPACT   [EXPAND] -> CO3
% ---------------------------------------------------------------------
\section{Societal, Health, Safety, Legal and Cultural Impact}
\label{sec:impacts}

The intended impact of this work is a decision-support tool small enough to
run where care is actually delivered, using the same contextual information a
clinician already has. The implications of such a tool span several
dimensions, and it is worth examining each in turn rather than treating
``impact'' as a single category.

\paragraph{Societal impact.}
The gap between the volume of skin lesions that need screening and the number
of dermatologists available to read them is the social problem this work
addresses. That gap is widest in low-resource settings: in countries such as
Bangladesh, specialist dermatological coverage is concentrated in major urban
centres, leaving rural populations with little access to early screening. A
model that is cheap to deploy and can run offline lowers the barrier to
first-line screening exactly where specialist coverage is thin, and so has the
potential to widen access rather than concentrate it further. Because the
system is assistive, its realistic social role is triage --- flagging lesions
that warrant referral --- rather than replacing the specialist judgement that
remains scarce.

\paragraph{Health and safety impact.}
The health benefit follows from the prognosis of skin cancer, which depends
heavily on early detection; melanoma in particular is highly treatable when
caught early. A tool that helps surface suspicious lesions sooner can, in
principle, improve outcomes. The safety considerations are the other side of
the same coin. The dangerous failure mode of a skin-lesion classifier is a
false negative on a malignant class, which is precisely why we optimise and
report balanced multi-class accuracy rather than overall accuracy: a model
that scored well on the common benign nevus class while missing the rare
malignant ones would be both numerically impressive and clinically unsafe. A
second safety concern is automation bias --- the tendency of users to defer to
a confident machine output. We mitigate this by designing the system to
produce a ranked set of possibilities rather than a single diagnosis, and by
positioning it explicitly as an aid to, not a replacement for, a clinician,
who remains in the loop for every decision.

\paragraph{Legal impact.}
Any clinical deployment of this system would sit within a medical-device
regulatory regime (Section~\ref{sec:standards}) and would require the
corresponding clearance before use on patients; the present work makes no
claim to such clearance. Liability for a missed or incorrect diagnosis, the
need for informed consent governing the use of patient data, and compliance
with data-protection law are all legal questions that a real deployment must
answer. By relying only on a public, de-identified research dataset and by
framing the output as decision support, this study stays within the bounds of
research use and leaves those clinical-deployment questions explicitly out of
scope.

\paragraph{Cultural impact.}
A specific and important limitation concerns skin-tone representation.
ISIC-2019 is assembled largely from European populations, and dermoscopic
collections in general over-represent fair skin. A model trained on such data
cannot be assumed to generalise to the darker and more diverse skin tones of a
South Asian population such as Bangladesh's, and deploying it there without
re-validation on representative local data would risk systematically worse
performance for the very patients it is meant to serve. We flag this as a
genuine equity concern rather than a solved problem, and note it again among
the limitations and future work (external validation on a local, representative
dataset). More broadly, the acceptance of AI-assisted diagnosis depends on
cultural attitudes toward automation in healthcare and on patient trust, which
reinforces the case for a transparent, assistive design rather than an opaque,
autonomous one.


% ---------------------------------------------------------------------
% 3.3  ENVIRONMENTAL IMPACT AND SUSTAINABILITY                   [NEW] -> CO4
% ---------------------------------------------------------------------
\section{Environmental Impact and Sustainability}
\label{sec:environmental-impact}

The environmental footprint of a deep-learning project has two distinct
components: the one-off cost of \emph{training} the models, and the recurring
cost of \emph{running} a deployed model over its lifetime. The efficiency
goal that drives this thesis is, in addition to a deployment argument, a
sustainability argument, and it is worth making that link explicit.

\paragraph{Training footprint.}
All experiments ran on a single workstation with an NVIDIA RTX~4070~Ti~Super
(a 285\,W board, of which we assume an average draw of roughly 0.25\,kW under
training load). The carbon cost of a single training run is then
\[
\text{CO}_2\text{eq} \;=\; P_{\text{avg}} \;\times\; t_{\text{train}} \;\times\; \text{EF},
\]
where $P_{\text{avg}}$ is the average board power, $t_{\text{train}}$ the
wall-clock training time, and $\text{EF}$ the grid emission factor. Bangladesh's
grid is fossil-dominated, with an official grid emission factor of
approximately $0.62\ \text{kg\,CO}_2/\text{kWh}$. Table~\ref{tab:carbon}
applies this to the two proposed models; the wall-clock times are the values
recorded in our training logs. Summed across the full study --- the five
baselines, both proposed models, and the ablation variants, on the order of a
dozen full training runs plus exploratory runs --- the total training energy is
on the order of tens of kWh, corresponding to roughly a few tens of kilograms
of CO\textsubscript{2}eq: a footprint comparable to a single short car
journey, and orders of magnitude below that of training the large ensembles
that top the ISIC-2019 leaderboard.

\begin{table}[h]
\centering
\caption{Estimated training energy and carbon for the two proposed models, at
an assumed average board power of 0.25\,kW and the Bangladesh grid emission
factor of $0.62\ \text{kg\,CO}_2/\text{kWh}$. Wall-clock times are illustrative
pending substitution with the exact values from the training logs.}
\label{tab:carbon}
\begin{tabular}{lccc}
\toprule
Model & Wall-clock (h) & Energy (kWh) & CO\textsubscript{2}eq (kg) \\
\midrule
Lightweight hybrid & \textbf{[REPLACE: $\approx$3]} & 0.75 & 0.47 \\
DEKAN (flagship)   & \textbf{[REPLACE: $\approx$9]} & 2.25 & 1.40 \\
\bottomrule
\end{tabular}
\end{table}

\paragraph{Inference footprint and sustainability.}
For a deployed clinical tool the recurring inference cost dominates the
lifetime footprint, and this is exactly where the design pays off. The
lightweight hybrid performs a single classification in 0.631 GMACs, so the
energy per inference is negligible, and the model is small enough to run on a
low-power edge device --- for example an NVIDIA Jetson Orin Nano, which runs
vision-transformer workloads within a power envelope of roughly 7--25\,W ---
rather than on a continuously powered server-class GPU drawing ten to twenty
times as much. A compact model also extends sustainability beyond energy: it
runs on inexpensive or already-available hardware, reducing the demand for new
high-end accelerators and the electronic waste that accompanies them, and it
operates without a permanent data-centre connection. By contrast, the
multi-resolution EfficientNet plus SENet-154 ensembles that define the
state of the art carry a far heavier per-inference cost. The two-tier design
of this thesis therefore offers an explicit efficiency-versus-accuracy choice
in which the lower-footprint option is also the more deployable one, aligning
the work with the goals of green and resource-efficient machine learning.

\paragraph{Honest accounting.}
We did not directly meter power consumption; the figures above are estimates
derived from the GPU's rated power and the national-average grid emission
factor, and the true values depend on the actual board utilisation and the
local energy mix. We report them as an order-of-magnitude account rather than
a precise measurement, which is sufficient to support the comparative claim
that the proposed models are low-footprint relative to leaderboard ensembles.


% ---------------------------------------------------------------------
% 3.4  ETHICAL ISSUES                                            [EXPAND] -> CO12
% ---------------------------------------------------------------------
\section{Ethical Issues}
\label{sec:ethics}

The ethical commitments of this work were stated briefly in the Ethics
Statement at the front of the thesis; here we set them out in the context of
the design.

\paragraph{Data ethics and privacy.}
The study uses only the publicly available ISIC-2019 dataset, released for
research by the International Skin Imaging Collaboration. No new human-subject
data were collected, and all images are de-identified at source. The patient
metadata --- age, sex, and anatomical site --- is used solely as model input
and is never used to re-identify a patient. The data is handled in accordance
with the ISIC data-use terms.

\paragraph{Fairness.}
Two design decisions are fairness decisions. First, the choice of balanced
multi-class accuracy as the primary metric forces the model to take the rare,
often malignant, classes as seriously as the common benign ones, rather than
optimising a figure the majority class would dominate. Second, missing
metadata is represented by a learned ``missing'' embedding rather than imputed,
so the model is never trained on a fabricated clinical value. We also
acknowledge openly (Section~\ref{sec:impacts}) the skin-tone representation
bias inherent in ISIC-2019, which is an unresolved fairness limitation rather
than a problem this work claims to have solved.

\paragraph{Responsible deployment and transparency.}
The system is designed to assist, not replace, a clinician: it outputs a
ranked set of possibilities and keeps a human in the loop, which guards
against the automation bias discussed earlier. In the same spirit of
transparency, we report our negative and inconclusive results honestly --- the
metadata benefit for the lightweight model could not be confirmed on a single
test-set evaluation, and the KAN classifier head did not improve over a plain
linear head --- rather than presenting only the favourable findings. Reporting
what did not work is itself an ethical obligation in a medical context, where
overstated performance can cause real harm.


% ---------------------------------------------------------------------
% 3.6  PROJECT MANAGEMENT PLAN                                   [NEW] -> CO10
% ---------------------------------------------------------------------
\section{Project Management Plan}
\label{sec:project-management}

The project was carried out by a team of five over three academic semesters,
from Spring~2025 to Summer~2026, under the supervision of
Dr.\ Md.\ Ashraful Alam. It was managed in the three phases that structure the
thesis programme --- Pre-Thesis~1, Pre-Thesis~2, and the Final Thesis --- with
the work broken down into the stages shown in
Table~\ref{tab:project-schedule}.

\paragraph{Work breakdown and schedule.}
The first phase established the problem and the foundations: the literature
review, the framing of the imbalance and efficiency problem, and the
construction of the data pipeline (perceptual-hash deduplication, the
lesion-grouped stratified split, colour constancy, and the memory-mapped
dataset). The second phase implemented and trained the five lightweight
baselines under a single shared recipe, and designed the lightweight hybrid.
The final phase designed and trained the higher-capacity DEKAN model, ran the
controlled ablations, carried out the full evaluation and analysis, and
produced the written thesis. Table~\ref{tab:project-schedule} maps these
stages onto the timeline.

\begin{table}[h]
\centering
\caption{Project schedule across the three thesis phases. A filled cell
indicates the stage was active in that phase.}
\label{tab:project-schedule}
\begin{tabular}{p{0.46\textwidth}ccc}
\toprule
Stage & Pre-Thesis 1 & Pre-Thesis 2 & Final Thesis \\
\midrule
Problem definition \& literature review        & $\blacksquare$ & $\square$ & $\square$ \\
Data acquisition \& preprocessing pipeline      & $\blacksquare$ & $\blacksquare$ & $\square$ \\
Baseline implementation (5 models)              & $\square$ & $\blacksquare$ & $\square$ \\
Lightweight hybrid design \& training           & $\square$ & $\blacksquare$ & $\blacksquare$ \\
DEKAN design \& training                        & $\square$ & $\square$ & $\blacksquare$ \\
Ablation studies                                & $\square$ & $\square$ & $\blacksquare$ \\
Evaluation, analysis \& figures                 & $\square$ & $\blacksquare$ & $\blacksquare$ \\
Thesis writing \& revision                      & $\blacksquare$ & $\blacksquare$ & $\blacksquare$ \\
\bottomrule
\end{tabular}
\end{table}

\paragraph{Roles and responsibilities.}
The work was divided across the team along the natural seams of the pipeline
--- data engineering and preprocessing, model architecture, training
infrastructure, evaluation and analysis, and writing and figures --- with
regular integration so that all models shared one identical training and
evaluation recipe. % [REPLACE: insert a one-line statement of which member led which area, e.g. "X led the data pipeline; Y and Z the model architectures; ...". This also supports the CO13 individual-contribution criterion.]
Coordination was maintained through version-controlled code and regular
supervisor meetings.

\paragraph{Tools and resources.}
The project used Git for version control of the modular codebase, the shared
single-GPU workstation described in Section~3.1.4 for all training, and the
configuration-driven pipeline that allows any result to be regenerated from a
saved specification. Keeping the entire study on one reproducible codebase and
one machine was itself a project-management decision: it removed the
environment drift that would otherwise make a five-person, multi-model
comparison unfair.

\paragraph{Milestones and deliverables.}
The principal milestones were the cleaned and split dataset, a working
baseline suite under the shared recipe, the budget-compliant lightweight
hybrid, the DEKAN flagship, the completed ablation matrix, and the final
written thesis together with the reproducibility appendix.


% ---------------------------------------------------------------------
% 3.7  RISK MANAGEMENT                                           [NEW] -> CO10
% ---------------------------------------------------------------------
\section{Risk Management}
\label{sec:risk-management}

A small number of technical risks were identified at the outset as the main
threats to a valid result, and the design was shaped to address them. Several
of the most consequential --- data leakage, rare-class collapse, and the
memory ceiling --- were mitigated by specific, deliberate engineering choices
that are described in Chapter~4; the risk register in
Table~\ref{tab:risk-register} records each risk, its likelihood and impact
before mitigation, the mitigation applied, and its residual status.

The leakage risk was the most insidious: because ISIC-2019 contains
near-duplicate photographs of the same lesion, naive splitting would have let
the same lesion appear in both training and test, inflating the reported
accuracy. Perceptual-hash deduplication followed by a lesion-grouped split
closed this off. The imbalance risk --- a model collapsing onto the common
classes and ignoring the rare malignant ones --- was countered with
Class-Balanced Focal Loss, a weighted sampler, and the choice of BMA as the
selection metric, and the training curves in Chapter~5 confirm no collapse
occurred. The 16\,GB memory ceiling and the Windows data-loader deadlock were
both engineering risks to feasibility, mitigated by mixed precision,
channels-last layout, gradient accumulation, and the memory-mapped,
zero-worker dataset. Two risks remain open and are carried forward as
limitations: all results come from a single seed, and the models have not been
validated on an external dataset, so variance and out-of-distribution
generalisation are not yet quantified.

\begin{table}[h]
\centering
\caption{Project risk register. L = likelihood, I = impact, before mitigation
(H/M/L).}
\label{tab:risk-register}
\begin{tabular}{p{0.26\textwidth}ccp{0.34\textwidth}c}
\toprule
Risk & L & I & Mitigation & Status \\
\midrule
Train/test leakage via duplicate lesion images & H & H & pHash deduplication + lesion-grouped, source-stratified split & Mitigated \\
Rare-class collapse under imbalance & H & H & Class-Balanced Focal Loss, weighted sampler, BMA selection metric & Mitigated \\
Overfitting of from-scratch transformer on $\sim$17k images & M & M & Pretrained CNN stem, stochastic depth, RandAugment, Mixup, weight EMA & Mitigated \\
16\,GB GPU memory ceiling & H & M & AMP, channels-last, gradient accumulation, memmap dataset & Mitigated \\
Windows data-loader deadlock & M & M & Zero-worker loading + memory-mapped array & Mitigated \\
Substantial metadata missingness & H & M & Learned ``missing'' embeddings, no imputation & Mitigated \\
Single-seed variance & H & M & Acknowledged; multi-seed reporting deferred & Open (future work) \\
No external-dataset validation & M & H & Acknowledged; external validation deferred & Open (future work) \\
KAN head failing to add value & M & L & Linear-head ablation retained and reported & Realised, handled \\
\bottomrule
\end{tabular}
\end{table}


% ---------------------------------------------------------------------
% 3.8  ECONOMIC ANALYSIS                                         [NEW] -> CO11
% ---------------------------------------------------------------------
\section{Economic Analysis}
\label{sec:economic-analysis}

The economic case for this work rests less on its development cost --- which is
modest and one-off --- than on the recurring cost of deployment, where a
compact model has a decisive advantage over both cloud-hosted inference and
the heavy ensembles that top the leaderboard.

\paragraph{Development and training cost.}
The development cost is dominated by the single workstation
(Section~3.1.4), a one-off capital expense of roughly USD 1{,}800--2{,}500,
together with the team's time over three semesters; the software stack is
entirely open-source and free. The compute cost of training is negligible: as
estimated in Section~\ref{sec:environmental-impact}, the whole study consumed
on the order of tens of kWh, which at a Bangladeshi electricity tariff of
roughly USD 0.08 per kWh amounts to only a few US dollars.

\paragraph{Deployment cost --- the decisive comparison.}
The lightweight hybrid, at 3.98M parameters and 0.631 GMACs, can be served
from an inexpensive edge device. An NVIDIA Jetson Orin Nano Super, which is
explicitly capable of running vision-transformer workloads, costs about
USD 249 as a one-off purchase and draws only a few watts, and can operate
offline --- a real advantage in rural clinics with unreliable connectivity.
The alternative, hosting the model on a cloud inference GPU, incurs a
recurring charge: even a modest T4-class instance costs roughly USD 0.15--0.81
per GPU-hour on demand, and a continuously available endpoint therefore costs
on the order of hundreds of US dollars per month, before accounting for data
egress and the privacy implications of sending patient images off-site.
Table~\ref{tab:tco} compares the two over a three-year horizon under a
plausible usage assumption; the edge option is roughly an order of magnitude
cheaper over the device's lifetime, and the gap only widens for the heavier
leaderboard ensembles, which require larger and more expensive accelerators to
serve at all.

\begin{table}[h]
\centering
\caption{Illustrative three-year total cost of ownership for serving the
lightweight model, edge versus cloud. Cloud assumes a T4-class instance at
USD 0.50/GPU-hour running 8 hours per day; edge assumes a one-off device plus
electricity at USD 0.08/kWh.}
\label{tab:tco}
\begin{tabular}{lccc}
\toprule
Option & Up-front (USD) & 3-year running (USD) & 3-year total (USD) \\
\midrule
Edge (Jetson Orin Nano Super, $\sim$10\,W) & 249 & $\sim$9   & $\sim$258 \\
Cloud (T4-class, 8\,h/day on demand)        & 0   & $\sim$4{,}380 & $\sim$4{,}380 \\
\bottomrule
\end{tabular}
\end{table}

\paragraph{Cost--benefit.}
Set against these costs, the benefit is earlier triage of suspicious lesions
in settings where specialist time is the scarce resource. The benefit is hard
to monetise precisely, but the asymmetry is clear: a single USD 249 device
amortised over many screenings per day yields an extremely low cost per
screening, and any malignant lesion caught earlier carries a large clinical
and economic value through avoided late-stage treatment. The accuracy gap
between the lightweight model (0.6081 BMA) and a leaderboard ensemble
($\approx$0.636) is small relative to the order-of-magnitude difference in
deployment cost and energy, which is the central economic argument for the
lightweight tier. Where a higher accuracy is worth a larger budget, DEKAN
(0.6438 BMA) provides it within the same design family, giving an institution
a principled choice between the two operating points rather than a single
fixed cost.
